\documentclass[11pt]{article}
\usepackage[a4paper,margin=1in]{geometry}
\usepackage{times}
\usepackage[hidelinks]{hyperref}
\usepackage{titlesec}
\usepackage{parskip}

\titleformat{\section}{\large\bfseries}{}{0em}{}
\titlespacing{\section}{0pt}{6pt}{2pt}

\begin{document}

\begin{center}
    {\LARGE \textbf{Mainul Islam}}\\[3pt]
    {\normalsize Statement of Research Interests}\\[2pt]
    {\small mainulislam@iut-dhaka.edu \ $\cdot$ \ +880-185-607-3265 \ $\cdot$ \
    \href{https://scholar.google.com/citations?user=mHfGHugAAAAJ&hl=en}{Google Scholar}}
\end{center}

\vspace{4pt}

My research interests lie at the intersection of \textbf{field robotics, physics-informed
machine learning, and computer vision}---building autonomous systems that perceive, reason,
and act reliably in the unstructured real world. Over the past three years I have worked across
the full autonomy stack, and I am now driven to deepen that work into rigorous, generalizable
research.

I am especially motivated by the gap between learning-based perception and the physical laws
that govern how robots actually move. In my undergraduate thesis, \textit{MonoSIM} (IEEE ICARA
2026), I built an open-source software-in-the-loop framework that couples monocular-vision lane
perception with model predictive and PID control for Ackermann vehicles---an early step toward
perception-to-control pipelines grounded in vehicle kinematics rather than data alone. This is
exactly the \emph{physics-informed} direction I want to pursue: embedding dynamical models,
kinematic constraints, and physical priors into vision and learning systems so that they remain
robust, sample-efficient, and trustworthy outside the training distribution.

My hands-on background gives this interest a practical anchor. I architected an end-to-end
\textbf{ROS~2} autonomy stack for a multi-robot ground-vehicle fleet---EKF sensor fusion,
space-time path planning, and real-time trajectory tracking---and developed a closed-form
inverse-kinematics solver and MoveIt~2 integration for a 7-DOF manipulator. In parallel, I have
published applied computer-vision and machine-learning research, including deep-learning
segmentation and fault diagnosis with explainable AI in \textit{Solar Energy} (IF~6.6) and
YOLO-based detection with vision-language OCR.

Through a PhD, I aim to unite these threads---\textbf{unmanned ground vehicles, manipulation,
and physics-informed perception}---into autonomous field robots that operate dependably in
demanding, real-world settings such as agricultural and off-road environments. The NSF-funded
work in your laboratory, spanning computer vision, physics-informed learning, and full-stack
field robotics, is the ideal setting to pursue this, and I would be excited to contribute my
end-to-end engineering experience to it.

\end{document}
